% Generated draft: requires mathematical and source review.
\documentclass[11pt]{article}
\usepackage[margin=25mm]{geometry}
\usepackage{fontspec}
\setmainfont{Latin Modern Roman}
\usepackage{amsmath,amssymb,mathtools,unicode-math}
\setmathfont{Latin Modern Math}
\usepackage{enumitem,needspace,xurl,hyperref}
\setlength{\parindent}{0pt}
\setlength{\parskip}{5pt}
\setlength{\emergencystretch}{3em}
\newcommand{\N}{\mathbb{N}}
\newcommand{\Z}{\mathbb{Z}}
\newcommand{\Q}{\mathbb{Q}}
\newcommand{\R}{\mathbb{R}}
\newcommand{\C}{\mathbb{C}}
\newcommand{\F}{\mathbb{F}}
\newcommand{\A}{\mathbb{A}}
\newcommand{\PP}{\mathbb{P}}
\newcommand{\E}{\mathbb{E}}
\newcommand{\eps}{\varepsilon}
\newcommand{\dd}{\,\mathrm{d}}
\begin{document}
\section*{UnsolvedMath ID: 800046}
\textbf{Source problem code:} catalogue entry
\subsection{Definitions and mathematical statement}
A semigroup is a nonempty set equipped with an associative binary operation. A finite semigroup can be presented by its finite multiplication table. Let \(\mathcal F\) be a finite set of finite semigroups, and write \(\mathcal V_{\mathrm{inc}}(\mathcal F)\) for the ``inclusive variety generated by \(\mathcal F\)'' in the terminology of the source.

The supplied record does not define an inclusive variety, an identical inclusion, or an elementary identical inclusion. Consequently, the following expresses the requested decision problem relative to those intended definitions; it does not supply the missing definitions. Let \(\mathcal E\) denote the intended collection of admissible inclusions, let \(\mathcal K\) denote the intended ambient class of semigroups, and let \(S\models\eta\) mean satisfaction of such an inclusion under the intended semantics. For \(B\subseteq\mathcal E\), put
\[
\operatorname{Mod}_{\mathcal K}(B)
=\{S\in\mathcal K: S\models\eta\text{ for every }\eta\in B\}.
\]
A finite basis for the generated inclusive variety means a finite set \(B\subseteq\mathcal E\) such that
\[
\mathcal V_{\mathrm{inc}}(\mathcal F)=\operatorname{Mod}_{\mathcal K}(B).
\]

\textbf{Problem (G. Mashevitzky).} Determine whether there exists an algorithm \(A\) which, on every input consisting of multiplication tables for a finite set \(\mathcal F\) of finite semigroups, halts and satisfies
\[
A(\mathcal F)=\mathrm{YES}
\quad\Longleftrightarrow\quad
\exists B\subseteq\mathcal E\;
\bigl(B\text{ is finite and }
\mathcal V_{\mathrm{inc}}(\mathcal F)=\operatorname{Mod}_{\mathcal K}(B)\bigr).
\]
Find such an algorithm or prove that none exists. The source writes ``(elementary) identical inclusions''; whether elementary inclusions are required, or two versions are intended, remains unresolved. The definitions of \(\mathcal E\), satisfaction, \(\mathcal K\), and the generated inclusive variety must be supplied before this becomes a fully specified decision problem. The source also does not specify whether the empty generating set is an allowed input.

The source compares this question with Tarski's problem for identities and reports that reference [M4] constructs a finite semigroup without a finite basis for its elementary identical inclusions. This reported example is background, not a resolution of the algorithmic question.
\subsection{Short English statement}
G. Mashevitzky asks for an algorithm deciding whether the inclusive variety generated by a finite set of finite semigroups has a finite basis of (elementary) identical inclusions, or a proof that no such algorithm exists. The supplied record omits the definitions needed to specify the inclusion language and its semantics precisely; the formulation above is conditional on those definitions.
\subsection{Sources}
\begin{itemize}
\item Local UnsolvedMath registry, permanent ID \(800046\), exact source code \texttt{catalogue entry}, title \emph{The Lyapin's Notebook --- Problem 48}. The supplied registry labels the problem open; this is a supplied status, not an independently verified mathematical assessment.
\item Bershadsky, Kublanovsky, Mashevitzky, \emph{The Lyapin's Notebook: Unsolved Problems in Semigroup Theory} (2026), Problem 48. Supplied source URL: \url{https://arxiv.org/abs/2604.04763}. The problem statement attributes the question to G. Mashevitzky. The registry lists Bershadsky, Kublanovsky, Mashevitzky as proposers and 2026 as the proposed year; the original proposal date is not established by the supplied excerpt.
\item Supplied provenance: accessed 2026-10-16; extraction method \texttt{arxiv\_latex}; upstream revision \texttt{d371ee917ac9647fbb06cd6fc529deec6f3dda82}. These are supplied dated provenance details, not an independent source verification.
\item Reference [M4], cited in the supplied statement for a finite semigroup having no finite basis for its elementary identical inclusions. No bibliographic details or separate URL for [M4] are supplied.
\item The comparison with Tarski's problem for identities is supplied by the statement; no separate reference or precise formulation of that comparison is provided.
\end{itemize}
\end{document}
